% !TEX root = ../notes.tex

\documentclass[../notes.tex]{subfiles}

\begin{document}

\section{March 16}
Today, we say more about Bousfield localization.

\subsection{Bousfield Localization}
Our next goal is to complete a spectrum. To this end, we will discuss the more general notion of Bousfield localization. We will need a size condition.
\begin{definition}[presentable]
	Fix a ring spectrum $R$. Then a full subcategory $\mc C\subseteq\mathrm{Mod}(R)$ is \textit{presentable} if and only if it satisfies the following.
	\begin{itemize}
		\item Closure: $\mc C$ is closed under de-suspensions and all colimits.
		\item Size: there is a set of objects generating $\mc C$ under colimits.
	\end{itemize}
\end{definition}
Here is the main theorem; it is an application of some Adjoint functor theorem.
\begin{theorem} \label{thm:bousfield-localization}
	Let $R$ be a ring spectrum, and let $\mc C\subseteq\mathrm{Mod}(R)$ be a presentable full subcategory. Then the inclusion $\mc C\to\mathrm{Mod}(R)$ preserves colimits and admits a right adjoint $G$ given by
	\[G_{\mc C}(M)\coloneqq\colim_{\substack{N\in\mc C\\f\colon N\to M}}N.\]
\end{theorem}
\begin{remark}
	We have tended to ignore set-theoretic conditions, but this one is important because we are trying to apply an Adjoint functor theorem. In particular, the colimit does not even really make sense if $\mc C$ is not even generated by a set of objects.
\end{remark}
\Cref{thm:bousfield-localization} gives rise to our localization.
\begin{definition}[Bousfield localization]
	Let $\mc C$ be a presentable full subcategory of $\mathrm{Mod}(R)$, where $R$ is a ring spectrum, and let $G_{\mc C}$ be the right adjoint of the inclusion. For an $R$-module $M$, the cofiber of the unit map $G_{\mc C}M\to M$ is the \textit{Bousfield localization}, labeled $L_{\mc C}M$. We will abbreviate the $\mc C$s in the sequel.
\end{definition}
\begin{remark}
	By construction, we have a fiber sequence
	\[G_{\mc C}(M)\to M\to L_{\mc C}(M).\]
\end{remark}
\begin{remark}
	One does not always require that $\mc C$ is closed under de-suspensions.
\end{remark}
It is useful to note that our ``localization'' does in fact produce a category of ``local'' objects.
\begin{definition}[local]
	Fix a ring $R$ and a presentable full subcategory $\mc C$ of $\mathrm{Mod}(R)$. Then an $R$-module $M$ is \textit{$\mc C$-local} if and only if every map $Y\to M$ is nullhomotopic for each $Y\in\mc C$.
\end{definition}
\begin{remark}
	Because we assumed that $\mc C$ is closed under de-suspensions, we see that showing that every map $Y\to M$ is nullhomotopic is equivalent to $\op{Maps}(Y,M)$ being nullhomotopic.
\end{remark}
\begin{remark} \label{rem:localization-is-local}
	Because $L_{\mc C}$ is a left adjoint, one sees that $L_{\mc C}M$ is $\mc C$-local. Indeed, one sees this by plugging into the adjunction.
\end{remark}
\begin{remark}
	It turns out that the collection of local objects is closed under shifts and limits (this is not hard to check directly), so the subcategory of local objects is stable.
\end{remark}
The moral is that our Bousfield localization comes with the additional data of ``local objects,'' and we have a localization functor which is left adjoint to the inclusion.

\subsection{Localization at a Module}
Here is a motivating example.
\begin{example}
	Let $\mc C$ denote the collection of $\mathbb S/p$-modules generated under all colimits from $\mathbb S/p$ and its shifts. A spectrum $X$ is then $\mc C$-local if and only if $\op{Hom}(\mathbb S/p,X)$ vanishes, which is equivalent to $X\otimes D(\mathbb S/p)$ vanishing. But $D(\mathbb S/p)$ is the fiber of $p\colon\mathbb S\to\mathbb S$, which admits $\mathbb S/p$ as a cofiber, so $D(\mathbb S/p)=\Sigma^{-1}\mathbb S/p$. Thus, $X$ is local if and only if $p\colon X\to X$ is an isomorphism. Now, $\mathbb S[1/p]$ can be defined to be the ring which admits these local objects as its module category, and it turns out that $L_{\mc C}=\mathbb S[1/p]\otimes_{\mathbb S}-$.
\end{example}
Let's generalize this construction.
\begin{definition}[acyclic]
	Fix an $R$-module $E$. Then a module $M$ is \textit{acyclic} if and only if $M\otimes_RE$ vanishes.
\end{definition}
We will need a set-theoretic construction of our acyclic modules. We will need to make a few more definitions.
\begin{proposition}
	Fix an $R$-module $E$. Then the category of acyclic modules is generated by a small subcategory.
\end{proposition}
\begin{proof}[Idea]
	In fact, the acyclic objects are generated by $G_E(E)$ and its shifts, which is done by a ``small object argument.''
\end{proof}
\begin{remark}
	Ravenel saw this proposition and his ``world shifted.'' He would then use this formalism to describe the homotopy groups of spheres, which is what we will soon explain.
\end{remark}
\begin{definition}[equivalence]
	Fix an $R$-module $E$. A map $f\colon M\to N$ of $R$-modules is an $E$-equiv\-alence if and only if $f\otimes E$ is an isomorphism.
\end{definition}
\begin{definition}[local]
	Fix an $R$-module $E$. An $R$-module $M$ is \textit{$E$-local} if and only if $\op{Hom}_R(Y,M)=0$ whenever $Y$ is acyclic. We will denote the corresponding localization functor by $L_E$.
\end{definition}
\begin{remark}
	It turns out that $E$-equivalences of $E$-local objects are isomorphisms.
\end{remark}
\begin{remark}
	If $E_1$ and $E_2$ are in the same Bousfield class, then we see that $E_1$ and $E_2$ have the same acyclic modules, so $L_{E_1}=L_{E_2}$.
\end{remark}
Here is a coherence check.
\begin{lemma}
	Fix an $\mathbb E_1$-$R$-algebra $E$. Then any $E$-module is $E$-local.
\end{lemma}
\begin{proof}
	Choose an $R$-module $M$ which admits an $E$-module structure. For any $Y$, we may compute
	\[\op{Hom}_R(Y,M)=\op{Hom}_E(Y\otimes_RE,M).\]
	Thus, if $Y$ is further acyclic, we see that this is $\op{Hom}_E(0,M)=0$.
\end{proof}
As one would expect from localization along some elements, one can give a ``universal property'' for our localization.
\begin{lemma} \label{lem:get-localization}
	Fix an $R$-module $E$. If $M$ is an $R$-module, then the map $M\to L_EM$ is characterized by requiring that
	\begin{listroman}
		\item $L_EM$ is local,
		\item and this map is an $E$-equivalence.
	\end{listroman}
\end{lemma}
\begin{proof}
	We will show the forward direction. Note $L_EM$ is local by \Cref{rem:localization-is-local}. As for (ii), we can hit the fiber sequence $G_EM\to M\to L_EM$ by $E$, so the latter is an $E$-equivalence if and only if $G_EM\otimes E=0$; this latter statement is true because $G_E$ outputs to acyclic objects.
\end{proof}
We are now able to complete.
\begin{example}
	Because $\mathbb S/p$ and $\mathbb S/p^4$ generate the same thick subcategory, they are in the same Bousfield class, so we get the same localization functor. Passing to $\mathbb S/p^4$ is convenient because $\mathbb S/p^4$ is an $\mathbb E_1$-ring, which is a theorem of Burkland. We now claim that
	\[L_{\mathbb S/p}X\stackrel?=\lim_nX/p^n.\]
	We use \Cref{lem:get-localization}: note that each $X/p^{4n}$ is a finite colimit of shifts of $X/p^4$, which is a module over $\mathbb S/p^4$ and hence local, so the target is local by stability. It remains to show that the natural map $X\to\lim_nX/p^n$ is an equivalence after hitting it with $-\otimes\mathbb S/p$. But the quotient in $(\lim_nX/p^n)/p$ can move inside the limit (limits commute with filtered colimits), so we receive $\lim_n(X/p^n/p)=\lim_n(X/p\oplus\Sigma X/p)$. By tracking the internal maps, we see that the limit is $X/p$!
\end{example}

\subsection{Idempotent Algebras}
In general, completion does not commute with all limits. To this end, we want the following definition. Recall that local objects are automatically closed under limits.
\begin{definition}[smashing]
	Fix a ring $R$. Then a localization of $R$-modules are \textit{smashing} if and only if the collection of local objects is closed under all colimits.
\end{definition}
Let's parameterize smashing localizations.
\begin{definition}[idempotent]
	An $R$-module $M$ is \textit{idempotent} if and only if the unit map $M\to M\otimes_RM$ is an equivalence.
\end{definition}
\begin{remark}
	An idempotent module admits a unique structure as an $\mathbb E_\infty$-$R$-algebra.
\end{remark}
\begin{remark}
	If $R$ is a Noetherian discrete ring, then the idempotent $R$-algebras are equivalent to the open subsets of $\Spec R$. This explains why idempotent algebras could be interesting to us.
\end{remark}
\begin{proposition}
	Fix a ring $R$. Then smashing localizations are exactly given by localization at the idempotent $R$-algebras.
\end{proposition}
\begin{proof}[Sketch]
	If $M$ is idempotent, then it turns out that $L_M\coloneqq M\otimes X$ is a smashing localization. Indeed, we will use \Cref{lem:get-localization} to recognize this localization: because $M$ is idempotent, the unit map $X\to M\otimes X$ is an $M$-equivalence, and the target is local because it is a module over the algebra structure on $M$.

	Conversely, suppose that $L$ is a smashing localization. Then we claim that $LX=LR\otimes_RX$, for which we use \Cref{lem:get-localization}. There is certainly a natural map
	\[X=R\otimes_RX\to LR\otimes_RX.\]
	The target is local because $X$ is a colimit of shifts of $R$, so the target is a colimit of shifts of $LR$, and $LR$ is local. By this expansion, one sees that this map is further an $LR$-equivalence because $L$ preserves colimits (because $L$ is smashing!).
\end{proof}
Thus, we are interested in producing some smashing localizations. Here is how to produce the category of acyclic objects.
\begin{proposition}
	Fix a full subcategory $\mc C$ of $\mathrm{Mod}(R)$ which is generated under colimits by a set of perfect $R$-modules. Then $L_\mc C$ is smashing.
\end{proposition}
\begin{proof}
	Observe that an $R$-module $M$ is local if and only if $\op{Mor}(Y,M)=0$ for all $Y\in\mc C$. We may equivalently check this for merely the $Y$ in a generating set of $\mc C$, for which we can take the perfect modules. Indeed, when perfect, we see that $\op{Mor}(Y,M)=DY\otimes_RM$, and the condition $DY\otimes_RM=0$ is now closed under arbitrary colimits!
\end{proof}
\begin{remark}
	In this situation, it turns out that $L_{\mc C}$ only depends on the thick subcategory generated by the generating perfect $R$-modules in $\mc C$.
\end{remark}
\begin{remark}
	It turns out that all smashing localizations arise from perfect $R$-modules in this way.
\end{remark}
\begin{example}
	Consider the discrete ring $R\coloneqq\FF_p\left[x^{1/p^\infty}\right]/(x)$. Then it turns out that idempotent algebras are still in bijection with smashing localizations, but they are not in bijection with open subsets of $\Spec R$. The problem is that this ring fails to be Noetherian. For example, $\lim\FF_p^{\otimes(\bullet+1)}$ turns out to be an idempotent algebra, which is the result of someone's thesis whose name was forgotten by Professor Hahn.
\end{example}
Anyway, let's write down the localizations that we are already have access to.
\begin{definition}
	Fix a prime $p$ and a nonnegative integer $n$. Then we let $L_n^f$ be the localization of the category of $p$-local spectra (namely, $\mathbb S_{(p)}$-modules) with acyclic objects $\mc C$ given by colimits given by those finite spectra of type at least $n+1$.
\end{definition}
\begin{remark}[Telescope]
	Ravenel has conjectured that all smashing localizations of $p$-local spectra take the form $L_n^f$ for some $n$. This was proven false in 2023.
\end{remark}
\begin{example}
	Note that $L_0^fM=L_0^f\mathbb S_{(p)}\otimes M=\mathbb S_{(p)}[1/p]\otimes M=\mathbb Q\otimes M$. For example, we see that $\pi_*L_0^f\mathbb S$ is concentrated in degree zero, where it is $\QQ$.
\end{example}
\begin{example}
	The homotopy groups $\pi_*L_1^f\mathbb S$ have been computed by Haynes Miller when $p>2$ and Mahowald when $p=2$. For example, when $p>2$, one has
	\[\pi_*L_1^f\mathbb S=\begin{cases}
		\ZZ_{(p)} & \text{if }n=0, \\
		\QQ_p/\ZZ_p & \text{if }n=-2, \\
		\ZZ/p^{k+1}\ZZ & \text{if }n+1=(p-1)p^km\text{ for some }k,m\text{ where }p\nmid m, \\
		0 & \text{otherwise}.
	\end{cases}\]
	Ravenel hoped that one could organize homotopy theory by complexity depending on which $L_n^f\mathbb S$ one would need (for finite $n$), and then one would be able to compute these localizations in finite time.
\end{example}
Let's explain how one could try to compute $\pi_*L_1^f\mathbb S$ by instead computing $\pi_*L_1^f(\mathbb S/p)$. One would eventually access $\pi_*L_1^f\mathbb S$ by trying to understand $\pi_*L_1^f(\mathbb S/p^\bullet)$ for larger and larger powers, pass to the completion, and then compare with the rationalization.
\begin{proposition}
	Suppose $X$ is a finite spectrum of type $n$ and admits a $v_n$-self map $f\colon\Sigma^kX\to X$. Then $L_n^fX=X\left[f^{-1}\right]$.
\end{proposition}
\begin{proof}
	We would like to recognize the natural map $X\to X\left[f^{-1}\right]$ as our localization map. The hard part is checking that $X\left[f^{-1}\right]$ is actually local. To this end, suppose that $Y$ is a finite spectrum of type greater than $n$, and we want to show that $\op{Hom}\left(Y,X\left[f^{-1}\right]\right)=0$. It is equivalent to check that $\id_{DY}\otimes f$ acts nilpotently on $DY\otimes X$, which is something we can check on Morava $K$-theories. We have two cases for our Morava $K$-theory $K(m)$.
	\begin{itemize}
		\item For $m\le n$, we see that $K(m)_*(DY\otimes X)=0$ by the K\"unneth formula (because $Y$ has type at least $n+1$).
		\item If $m>n$, then the action is still nilpotent because $f$ is a $v_n$-self map!
	\end{itemize}
	We will omit checking that we vanish under $(\FF_p)_*$ as well, which should hold by some limiting argument.

	It remains to check that $X\to X\left[f^{-1}\right]$ is an $L_n^f$-equivalence. It is enough to check that the fiber is a colimit of objects which are killed by $L_n^f$. Well, the cofiber of $X\to X\left[f^{-1}\right]$, but this is the colimit of
	\[0\to X/f\to X/f^2\to\cdots.\]
	Now, each $X/f^n$ is of type $n+1$ (it is the cofiber of a $v_n$-self map), which are acyclic for $L_n^f$ by definition of the localization.
\end{proof}
\begin{example}
	Thus, we see that
	\[L_1^f\mathbb S/p=\mathbb S/p\left[v_1^{-1}\right].\]
	One idea to compute these homotopy groups would be to invert $v_1$ on the $E_2$-page of the Adams--Novikov spectral sequence for $\mathbb S/p$; this does give the correct answer, which is $\FF_p[v_1,1/v_1]\otimes\Lambda(\zeta)$, where $\zeta$ has degree $-1$, and $v_1$ has degree $2p-2$. Do note that this idea does not really make sense immediately because inverting $v_1$ on the $E_2$-page does not even have to produce a convergent spectral sequence. Ravenel showed that this idea would work for $p$ large, conditional on his Telescope conjecture, which we recall is false: it turns out to fail for $L_2^f$!
\end{example}

\end{document}